\begin{multicols}{2}
   \topic{Paging and Demand Paging}
   \subtopic{Page Fault}

\begin{itemize}[noitemsep]
    \item Present bit in page table entry: indicates if a page of a process resides in memory or not
    \item When translating VA to PA, MMU reads present bit
    \item If page present in memory, directly accessed
   \item If page not in memory, MMU raises a trap to the OS - page fault
\end{itemize}


\subtopic{Translation Lookaside Buffer (TLB)}
\begin{itemize}[noitemsep]
    \item A cache of recent VA-PA mappings
    \item To translate VA to PA, MMU first looks up TLB
    \item If TLB hit, PA can be directly used
    \item If TLB miss, then MMU performs additional memory accesses to ``walk'' page table
    \item TLB misses are expensive (multiple memory accesses)
    \begin{itemize}[noitemsep]
        \item Locality of reference helps to have high hit rate
    \end{itemize}
    \item TLB entries may become invalid on context switch and change of page tables
\end{itemize}

\subtopic{Demand Paging}
\begin{itemize}[noitemsep]
    \item Memory management technique where pages are loaded into physical memory only when needed (on first access).
    \item Reduces memory usage by avoiding loading unused pages.
    \item If a required page is not in memory, a page fault occurs $\to$ OS loads the page from disk into a free frame.
    \item Common in virtual memory systems for efficient resource use.
\end{itemize}

\columnbreak

Page algos
\end{multicols}